\chapter{Evaluation Methodology}
\label{chap:evaluation}

\section{Experimental setup}
If the assignment is a literature review, use this chapter to compare the
evaluation methods in the papers. If you run experiments, record at least:
\begin{itemize}
  \item CPU model, cache hierarchy, operating system, compiler version, and flags;
  \item input sizes, storage layout, initialisation, and correctness criteria;
  \item repetitions, summary statistics, and warm-up policy;
  \item timing boundaries and the exact meaning of any hardware counters.
\end{itemize}
Keep allocation and initialisation outside the timed region unless they are
part of the operation being evaluated. Explain how you prevent the compiler
from removing work whose output is never used.

\section{Presenting results}
For a baseline time $t_b$ and an optimised time $t_o$, define speedup as
\begin{equation}
  S = \frac{t_b}{t_o}, \qquad t_o>0.
  \label{eq:speedup}
\end{equation}
Values above one indicate a shorter execution time. \Cref{tab:results}
uses deliberately synthetic values to illustrate units, alignment, and captions.
These values are not measurements or results from the cited papers.

\begin{table}[htbp]
  \centering
  \caption{Formatting example only: synthetic execution times and speedups.}
  \label{tab:results}
  \begin{tabular}{lrrr}
    \toprule
    Variant & Time (ms) & Speedup & Repetitions \\
    \midrule
    Baseline & 120.0 & $1.00\times$ & 10 \\
    Interchanged & 80.0 & $1.50\times$ & 10 \\
    Tiled & 60.0 & $2.00\times$ & 10 \\
    \bottomrule
  \end{tabular}
\end{table}

\section{Threats to validity}
\todo{Discuss measurement variability, input selection, compiler-generated
transformations, cache warm-up, and whether the hardware represents the setting
you want to study. For a literature review, explain differences in experimental
conditions across papers.}
